\chapter{Brittle Nails (Nail Changes)}

\section*{Definition}
Nails that are weak, thin, split, peel, or break easily, reflecting impaired keratinisation or nutritional deficiencies affecting the nail matrix and plate.

\section*{When to See a Doctor}

\begin{itemize}
 \item Brittle nails with nail discolouration, thickening, separation from nail bed, or associated with fatigue, pallor, or weight changes
\end{itemize}

\section*{How It Develops}
Nail plate formation occurs in the nail matrix through keratinocyte differentiation and keratinisation. Oestrogen receptors in the nail matrix regulate keratin production; declining oestrogen during menopause reduces nail thickness and integrity. Reduced oil production (sebum) also affects the nail cuticle and hydration. Nutritional deficiencies, repeated wet-dry cycles (housework, swimming), and trauma from manicure techniques further compromise nail structure.

\section*{Causes}
\begin{itemize}
 \item Hormonal changes of perimenopause and menopause
 \item Iron deficiency anaemia
 \item Zinc deficiency
 \item Biotin deficiency
 \item Hypothyroidism or hyperthyroidism
 \item Repeated exposure to water and detergents
 \item Fungal nail infection (onychomycosis)
 \item Psoriasis (nail pitting, onycholysis)
 \item Trauma or aggressive manicure techniques
 \item Raynaud phenomenon or peripheral vascular disease
\end{itemize}

\section*{Related Symptoms}
\begin{itemize}
 \item Hair loss
 \item Dry skin
 \item Fatigue
 \item Anaemia
 \item Hypothyroidism
\end{itemize}


\section*{Medical Evaluation}
\begin{itemize}
 \item Your doctor will examine the nails for pitting, onycholysis, discolouration, thickening, and signs of fungal infection or psoriasis.
 \item Tests are guided by the history and examination. A blood count, ferritin, or thyroid testing may be considered when symptoms suggest anaemia, iron deficiency, or thyroid disease; zinc and biotin testing is not routine.
 \item Nail clippings for fungal culture or periodic acid-Schiff (PAS) staining are obtained if onychomycosis is suspected.
 \item Referral to a dermatologist is arranged for nail psoriasis, persistent fungal infection despite topical treatment, or when the diagnosis remains unclear.
\end{itemize}

\section*{Treatment Options}
\begin{itemize}
 \item Biotin supplements have not been shown to reliably improve brittle nails when there is no deficiency. High-dose biotin can interfere with some laboratory tests, so tell clinicians about its use and correct a confirmed deficiency with professional guidance.
 \item Suspected fungal infection should be confirmed when practical before prolonged treatment. Topical or oral antifungals are selected according to the organism, extent, medicines, and liver health.
 \item Iron replacement therapy is prescribed when iron deficiency anaemia is identified as the underlying cause.
 \item Topical corticosteroid or vitamin D analogue creams are used for nail psoriasis, with systemic therapy reserved for severe or extensive disease.
\end{itemize}

\section*{Self-Care and Home Management}
\begin{itemize}
 \item Keep nails short to reduce breakage and file in one direction (not back and forth) to prevent splitting.
 \item Apply a moisturising cream or oil (jojoba oil, shea butter) to the nail plate and cuticles after hand washing and before bed.
 \item Wear gloves when washing dishes, cleaning, or gardening to protect nails from water and chemical exposure.
 \item Avoid using nail polish removers containing acetone, which dries out the nail plate and worsens brittleness.
\end{itemize}

\section*{References}
\begin{thebibliography}{99}
\bibitem{brittle_nails:ref1} Iorizzo M, Piraccini BM, Tosti A. ``Brittle nails.'' J Cosmet Dermatol. 2004;3(3):138-144.
\bibitem{brittle_nails:ref2} van de Kerkhof PCM, Pasch MC, Scher RK, et al. ``Nail psoriasis and nail lichen planus.'' In: Scher RK, Daniel CR, eds. Nails: Diagnosis, Therapy, Surgery. 3rd ed. Saunders; 2005.
\bibitem{brittle_nails:ref3} Holzberg M. ``Nail disorders in the elderly.'' Dermatol Clin. 2021;39(2):275-286.
\bibitem{brittle_nails:ref4} de Berker D. ``Nail anatomy and physiology.'' In: Fitzpatrick's Dermatology. 9th ed. McGraw-Hill; 2019.
\bibitem{brittle_nails:ref5} Singal A, Arora R. ``Nail as a window of systemic diseases.'' Indian Dermatol Online J. 2015;6(2):67-74.
\bibitem{brittle_nails:ref6} UpToDate. ``Nail disorders: Brittle nails and onychoschizia.'' Wolters Kluwer; 2025.
\end{thebibliography}
